% GENERATED FILE. Edit canonical lessons, then run npm run generate:books.
% canonical-source-hash: fnv1a64:92c0a6d9f98106dc
% canonical-lessons: KA-C145-kahi, KA-C145-huli2, KA-C145-sihi, KA-C145-somari, KA-C145-jana

\chapter{Bitter, Sour, Sweet, Lazy, Clever}
\label{ch:ka-bitter-sour-sweet-lazy-clever}
\hlchaptermodality{145}

\begin{chapteropening}
\textbf{By the end of this chapter:} \emph{I can say bitter, sour, sweet, lazy and clever.}

Complete the last lesson of chapter 145: i can say bitter, sour, sweet, lazy and clever.
\end{chapteropening}

\section[kahi]{\kn{ಕಹಿ} (kahi) — bitter}

\label{lesson:KA-C145-kahi}



\begin{quote}
Before the new one: say the Kannada for grey, then the Kannada for spicy.
\end{quote}

\subsection*{You'll want to know: \kn{ಕಹಿ}}
\textbf{\kn{ಕಹಿ}} — \emph{kahi} — \textquotedblleft{}bitter\textquotedblright{}.

\begin{grammarlens}[title={What you've built}]
One more word for this chapter.
\end{grammarlens}

\subsection*{Your turn}
\begin{itemize}
\raggedright
  \item \emph{Say it:} \emph{kahi}
  \item \emph{Say it:} \emph{kahi}, once more
  \item \emph{Say it:} say \emph{kahi}
  \item \emph{Recall:} say the Kannada for grey, then the Kannada for spicy, then say \emph{kahi} again
\end{itemize}

\begin{tcolorbox}[breakable,colback=teal!4,colframe=teal!35!black,title={Before you move on}]
What does \kn{ಕಹಿ} mean? (\textquotedblleft{}Bitter\textquotedblright{}.) Say it once more.
\end{tcolorbox}
\section[huḷi]{\kn{ಹುಳಿ} (huḷi) — sour}

\label{lesson:KA-C145-huli2}



\begin{quote}
Before the new one: say the Kannada for spicy, then the Kannada for bitter.
\end{quote}

\subsection*{You'll want to know: \kn{ಹುಳಿ}}
\textbf{\kn{ಹುಳಿ}} — \emph{huḷi} — \textquotedblleft{}sour\textquotedblright{}.

\begin{grammarlens}[title={What you've built}]
One more word for this chapter.
\end{grammarlens}

\subsection*{Your turn}
\begin{itemize}
\raggedright
  \item \emph{Say it:} \emph{huḷi}
  \item \emph{Say it:} \emph{huḷi}, once more
  \item \emph{Say it:} say \emph{huḷi}
  \item \emph{Recall:} say the Kannada for spicy, then the Kannada for bitter, then say \emph{huḷi} again
\end{itemize}

\begin{tcolorbox}[breakable,colback=teal!4,colframe=teal!35!black,title={Before you move on}]
What does \kn{ಹುಳಿ} mean? (\textquotedblleft{}Sour\textquotedblright{}.) Say it once more.
\end{tcolorbox}
\section[sihi]{\kn{ಸಿಹಿ} (sihi) — sweet}

\label{lesson:KA-C145-sihi}



\begin{quote}
Before the new one: say the Kannada for bitter, then the Kannada for sour.
\end{quote}

\subsection*{You'll want to know: \kn{ಸಿಹಿ}}
\textbf{\kn{ಸಿಹಿ}} — \emph{sihi} — \textquotedblleft{}sweet\textquotedblright{}.

\begin{grammarlens}[title={What you've built}]
One more word for this chapter.
\end{grammarlens}

\subsection*{Your turn}
\begin{itemize}
\raggedright
  \item \emph{Say it:} \emph{sihi}
  \item \emph{Say it:} \emph{sihi}, once more
  \item \emph{Say it:} say \emph{sihi}
  \item \emph{Recall:} say the Kannada for bitter, then the Kannada for sour, then say \emph{sihi} again
\end{itemize}

\begin{tcolorbox}[breakable,colback=teal!4,colframe=teal!35!black,title={Before you move on}]
What does \kn{ಸಿಹಿ} mean? (\textquotedblleft{}Sweet\textquotedblright{}.) Say it once more.
\end{tcolorbox}
\section[sōmāri]{\kn{ಸೋಮಾರಿ} (sōmāri) — lazy}

\label{lesson:KA-C145-somari}



\begin{quote}
Before the new one: say the Kannada for sour, then the Kannada for sweet.
\end{quote}

\subsection*{You'll want to know: \kn{ಸೋಮಾರಿ}}
\textbf{\kn{ಸೋಮಾರಿ}} — \emph{sōmāri} — \textquotedblleft{}lazy\textquotedblright{}.

\begin{grammarlens}[title={What you've built}]
One more word for this chapter.
\end{grammarlens}

\subsection*{Your turn}
\begin{itemize}
\raggedright
  \item \emph{Say it:} \emph{sōmāri}
  \item \emph{Say it:} \emph{sōmāri}, once more
  \item \emph{Say it:} say \emph{sōmāri}
  \item \emph{Recall:} say the Kannada for sour, then the Kannada for sweet, then say \emph{sōmāri} again
\end{itemize}

\begin{tcolorbox}[breakable,colback=teal!4,colframe=teal!35!black,title={Before you move on}]
What does \kn{ಸೋಮಾರಿ} mean? (\textquotedblleft{}Lazy\textquotedblright{}.) Say it once more.
\end{tcolorbox}
\section[jāṇa]{\kn{ಜಾಣ} (jāṇa) — clever}

\label{lesson:KA-C145-jana}



\begin{quote}
Before the new one: say the Kannada for sweet, then the Kannada for lazy.
\end{quote}

\subsection*{You'll want to know: \kn{ಜಾಣ}}
\textbf{\kn{ಜಾಣ}} — \emph{jāṇa} — \textquotedblleft{}clever\textquotedblright{}.

\begin{grammarlens}[title={What you've built}]
One more word for this chapter.
\end{grammarlens}

\subsection*{Your turn}
\begin{itemize}
\raggedright
  \item \emph{Say it:} \emph{jāṇa}
  \item \emph{Say it:} \emph{jāṇa}, once more
  \item \emph{Say it:} say \emph{jāṇa}
  \item \emph{Recall:} say the Kannada for sweet, then the Kannada for lazy, then say \emph{jāṇa} again
\end{itemize}

\begin{tcolorbox}[breakable,colback=teal!4,colframe=teal!35!black,title={Before you move on}]
What does \kn{ಜಾಣ} mean? (\textquotedblleft{}Clever\textquotedblright{}.) Say it once more.
\end{tcolorbox}
